\documentclass{article}
\usepackage{neurips_2025}

\usepackage[utf8]{inputenc}
\usepackage[T1]{fontenc}
\usepackage{hyperref}
\usepackage{url}
\usepackage{booktabs}
\usepackage{amsfonts}
\usepackage{nicefrac}
\usepackage{microtype}
\usepackage{xcolor}
\usepackage{graphicx}
\usepackage{float}

\title{Détection d’animaux camouflés par information temporelle et modèles Vision-Langage}

\begin{document}
\maketitle

\section{Résultats et Conclusions préliminaires}

\subsection{Objectifs}

\textbf{V1 et V2 (CNN).} 
L'objectif de V1 et V2 est d’évaluer la capacité d’un CNN pré-entraîné (ResNet) pour la détection d’animaux camouflés.

\textbf{V3 (hybride: mouvement + CLIP).}
L’objectif de V3 est de proposer une méthode légère pour détecter des animaux camouflés sans entraînement supervisé lourd. L’idée est de combiner une information temporelle simple, extraite entre images consécutives, avec un modèle Vision-Language capable de sélectionner la région la plus susceptible de contenir un animal.

\textbf{V4 (OWL-ViT).} 
L’objectif de V4 est d’évaluer une approche Vision-Language pour la détection d’animaux camouflés sur MoCA, en combinant OWL-ViT (prompts textuels) avec un raffinement temporel léger, puis en comparant une version zéro-shot et une version fine-tunée.

% \paragraph{Méthode}
% Pour V1 et V2, les étapes sont très similaires. Pour V1, on intègre à l'ensemble d'entrainement 29 photos par animal et une photo pour l'ensemble de validation. Pour V2, on intègre à l'ensemble d'entrainement 30 photos par animal et 10 photos par animal pour l'ensemble de validation. Le modèle pour V2 applique un vote de majorité pour trancher sur l'animal détecté.

\subsection{Méthodes}

\textbf{V1 et V2.}
Les deux approches sont similaires. Pour V1, 29 images par animal sont utilisées pour l’entraînement et 1 pour la validation. Pour V2, 30 images sont utilisées pour l’entraînement et 10 pour la validation. Le modèle V2 applique un vote de majorité pour déterminer l’animal détecté.

\textbf{V3.}
L’approche suit quatre étapes. D’abord, le mouvement est estimé entre deux frames successives afin de faire émerger des régions candidates, ensuite, chaque région est recadrée, puis comparée avec CLIP à une série de prompts décrivant un animal camouflé ou caché, enfin, un score hybride combinant mouvement et CLIP est utilisé pour ne conserver qu’un seul candidat finale par image. Plus précisément, le score final est calculé sous la forme : $S = \alpha \cdot S_{\text{motion}} + \beta \cdot S_{\text{CLIP}}$
% \[
% S = \alpha \cdot S_{\text{motion}} + \beta \cdot S_{\text{CLIP}},
% \]
où $\alpha$ pondère l’importance du mouvement et $\beta$ celle de la similarité sémantique fournie par CLIP. Dans toutes les configurations testées, une seule détection finale est gardée par frame.

\textbf{V4.}
L’approche repose sur une chaîne complète : détection frame-level avec OWL-ViT guidé par des prompts textuels, filtrage des prédictions (top-$k$, seuil de score), puis raffinement et agrégation temporelle. L’évaluation est réalisée sur cinq vidéos MoCA avec AP@0.5, précision, rappel et F1. Le fine-tuning est effectué sur les frames annotées avec une séparation train/validation par vidéo. Plusieurs configurations (prompts, taux d’apprentissage, poids de loss) sont testées en validation croisée leave-one-video-out. Le meilleur modèle est sélectionné selon le F1, puis un ajustement du seuil de score est appliqué pour retenir la configuration finale.

% \subsection{Métriques utilisées}
% Pour cette partie, 3 métriques ont été utilisées.

% \begin{itemize}
%    \item \textbf{Accuracy} $\left( \frac{TP+TN}{TP+FP+TN+FN} \right)$ : proportion de bonnes prédictions. La métrique la plus simple.

%     \item \textbf{Précision} $\left( \frac{TP}{TP + FP} \right)$ : proportion de détections prédites qui sont correctes. Elle pénalise surtout les faux positifs.

%     \item \textbf{Rappel} $\left( \frac{TP}{TP + FN} \right)$ : proportion d’objets annotés effectivement retrouvés. Il pénalise surtout les objets manqués.

\subsection{Métriques}

Nous utilisons des métriques standard pour évaluer les performances de classification et de détection, afin de conserver une comparaison cohérente entre les approches.

\begin{itemize}
    \item \textbf{Accuracy} $\left(\frac{TP+TN}{TP+FP+TN+FN}\right)$ : proportion de prédictions correctes.

    \item \textbf{Précision} $\left(\frac{TP}{TP+FP}\right)$ : proportion de prédictions positives correctes; pénalise les faux positifs.

    \item \textbf{Rappel} $\left(\frac{TP}{TP+FN}\right)$ : proportion d’objets détectés parmi les objets annotés; pénalise les faux négatifs.

    \item \textbf{F1-score} $\left(\frac{2PR}{P+R}\right)$ : moyenne harmonique entre précision et rappel.

    \item \textbf{IoU} (Intersection over Union) : mesure le chevauchement entre la boîte prédite et la boîte annotée; une détection est correcte si $\text{IoU} \geq 0.5$.

    \item \textbf{AP@0.5} : aire sous la courbe précision-rappel en considérant $\text{IoU} \geq 0.5$, synthétisant la performance globale sur différents seuils.
\end{itemize}

Les scores macro correspondent à la moyenne des performances sur les 5 vidéos évaluées.


\subsection{Résultats}

\begin{table}[h]
\centering
\begin{tabular}{lcccc}
\hline
\textbf{Version} & \textbf{Accuracy} & \textbf{Précision} & \textbf{Rappel} & \textbf{F1} \\
& (macro) & (macro) & (macro)& (macro)\\
\hline
Resnet static & 0.481 & 0.431 & 0.481 & 0.441 \\
Resnet temporelle & 0.556 & 0.459 & 0.556 & 0.481 \\
\hline
\end{tabular}
\caption{Performances des modèles Resnet static vs temporelle.}
\end{table}

\begin{table}[H]
\centering
\begin{tabular}{lcccc}
\hline
\textbf{Configuration V3} & \textbf{AP@0.5} & \textbf{Précision} & \textbf{Rappel} & \textbf{F1} \\
 & \textbf{(macro)} & \textbf{(macro)} & \textbf{(macro)} & \textbf{(macro)} \\
\hline
Config. A & 0.0332 & 0.0772 & 0.0752 & 0.0762 \\
Config. C & 0.0127 & 0.0401 & 0.0391 & 0.0396 \\
Config. E & 0.0402 & 0.0662 & 0.0627 & 0.0641 \\
\hline
\end{tabular}
\caption{Résultats macro de la configuration E sur 16 vidéos de développement.}
\label{tab:v3_results}
\end{table}

\begin{table}[H]
\centering
\begin{tabular}{lccccc}
\hline
\textbf{Configuration V4} & \textbf{Seuil} & \textbf{AP@0.5} & \textbf{Précision} & \textbf{Rappel} & \textbf{F1} \\
 & \textbf{score} & \textbf{(macro)} & \textbf{(macro)} & \textbf{(macro)} & \textbf{(macro)} \\
\hline
Baseline zéro-shot & 0.10 & 0.7574 & 0.9196 & 0.7857 & 0.8365 \\
Fine-tuning & 0.08 & 0.8176 & 0.9199 & 0.7978 & 0.8438 \\
\hline
\end{tabular}
\caption{Résultats V4 sur 5 vidéos MoCA.}
\end{table}

\subsection{Analyse}
\textbf{V1 et V2.}
Les différentes métriques donnent toutes des très bons résultats. Le vote de majorité pour V2 semble donner un résultat légèrement à V1 pour une mesure de précision, ce qui est étonnant vu que le V2 devrait avoir des résultats avec une plus faible variance que V1, mais les résultats restent très hauts. 

\textbf{V3.}
La configuration E, qui renforçait davantage la composante CLIP, resserrait les boîtes proposées et désactivait le fallback plein écran, n’a pas permis d’améliorer les performances globales. Les résultats macro obtenus sur les 16 vidéos de développement sont présentés au tableau~\ref{tab:v3_results}. Ces résultats montrent que le simple fait d’augmenter le poids de CLIP ou d’utiliser un backbone plus grand ne va nécessairement pas compenser une mauvaise qualité des propositions initiales. (il reste pas mal d'autres paramètres à explorer ou d'autre tweak pour améliorer les performances de cet approche)

\textbf{V4.}
Le fine-tuning améliore modérément les performances: AP@0.5 passe de 0.7574 à 0.8176 (+0.0602) et F1 de 0.8365 à 0.8438 (+0.0073), avec une précision stable (0.9196→0.9199) et un léger gain en rappel (0.7857→0.7978). L’amélioration principale concerne donc le score global de détection.

\subsection{Limites observées}
\textbf{V1 et V2.}
Le modèle ne détecte pas nécessairement "l'animal". Le CNN est entrainé sur toute l'image, donc dû à l'incapacité d'interprétation des CNNs, on ne peut être sûre que le modèle détecte exactement l'animal. Il est fort possible que le modèle utilise l'environnement pour reconnaitre l'animal, ce qui est contre la tâche apparente.  On a là la grande limitation des modèles CNN pour cette tâche.

\textbf{V3.}
Les résultats obtenus montrent que, dans sa forme actuelle, l’approche V3 reste limitée en performance sur l’ensemble de développement. Cela suggère que la combinaison entre propositions de mouvement et reranking par CLIP, telle qu’implantée ici, ne suffit pas encore à produire des détections robustes.

\textbf{V4.}
Les performances restent variables selon les vidéos, notamment en cas de camouflage fort (boîtes imprécises ou objets manqués). De plus, les résultats dépendent du choix des seuils (IoU, score), nécessitant un protocole d’évaluation cohérent.

\subsection{Conclusion}
\textbf{V1 et V2.}
Vu les résultats d'exactitude de V1 et V2, nous pouvons nous couvaincre que les modèles CNNs sont des bonnes options pour l'introduction de la détection d'animaux camouflés, mais la limitation des CNNs évoqués plus hauts rendent les modèles moins pertinents pour la tâche à faire et pousse à la recherche de meilleurs modèles.

\textbf{V3.}
V3 constitue une approche hybride simple, modulaire et interprétable, fondée sur la complémentarité entre information temporelle et information sémantique. Toutefois, dans sa forme actuelle, ses performances restent très limitées.

\textbf{V4.}
À ce stade, V4 est opérationnel et montre une amélioration mesurable après fine-tuning et calibration, sans rupture majeure par rapport à la baseline. La suite du travail pour V4 sera de consolider cette version dans l’évaluation finale commune avec les autres variantes du projet.

% \newpage

\section*{Annexe}
\subsection*{Approche V3: hybride mouvement + CLIP}

% Plus précisément, le score final est calculé sous la forme :
% \[
% S = \alpha \cdot S_{\text{motion}} + \beta \cdot S_{\text{CLIP}},
% \]
% où $\alpha$ pondère l’importance du mouvement et $\beta$ celle de la similarité sémantique fournie par CLIP. Dans toutes les configurations testées, une seule détection finale est gardée par frame.

\paragraph{Configurations explorées}
Pour garder cette section concise, nous avons retenu trois configurations représentatives : A, C et E. Les principaux paramètres étudiés sont le seuil de différence (\texttt{DIFF\_THRESHOLD}), la taille minimale des régions (\texttt{MIN\_AREA}), l’expansion des boîtes (\texttt{BOX\_EXPAND}), le nombre maximal de propositions (\texttt{PROPOSAL\_TOP\_K}), le backbone CLIP, ainsi que les poids $\alpha$ et $\beta$.

\begin{table}[H]
\centering
\begin{tabular}{lccccccc}
\hline
\textbf{Config.} & \textbf{Diff.} & \textbf{Min area} & \textbf{Expand} & \textbf{Top-K} & \textbf{CLIP} & \boldmath$\alpha$ & \boldmath$\beta$ \\
\hline
A & 25 & 400 & 0.15 & 10 & ViT-B-32 & 0.35 & 0.65 \\
C & 14 & 150 & 0.12 & 10 & ViT-B-32 & 0.15 & 0.85 \\
E & 12 & 150 & 0.05 & 15 & ViT-L-14 & 0.30 & 0.70 \\
\hline
\end{tabular}
\caption{Résumé des principales configurations V3 testées.}
\end{table}
\paragraph{Pistes d’amélioration}
Une amélioration de la méthode serait peut-être de ne plus comparer le frame $t$ au  frame $t-1$, mais potentiellement une méthode qui pourra améliorer les score c'est de prendre le frame $t$ avec le frame $t-k$, avec un pas temporel $k$ (à voir !). Cette modification va laisser le temps aux animaux lents pour se déplacer physiquement dans la scène, ce qui pourrait produire un masque de mouvement plus grand, plus net et donc plus exploitable pour la génération de propositions.

\end{document}
